% Extracted Guiding Questions
\documentclass[12pt]{article}
\begin{document}
\title{Global Warming Science: guiding questions for weekly essays}\author{Eli Tziperman}\date{}\maketitle
\section*{Greenhouse}
  \textbf{Guiding questions to be addressed in your report:}
    \begin{enumerate}
    \item What does it mean that the atmosphere is already saturated with respect to LW absorption at the CO$_2$ absorption bands? Why is temperature still increasing in response to CO$_2$ increase?
    \item What is the radiative forcing of CO$_2$, and how is it calculated?
    \item Should we worry about CO$_2$ increase given that H$_2$O is a much more powerful greenhouse gas?
    \item What are the sources of uncertainty in a calculation of radiative response to CO$_2$ increase by the year 2100?
    \item How is the temperature expected to increase as the CO$_2$ doubles once? Twice?
    \item How do we quantify the effects of other greenhouse gases, and how would you use that to set a policy regarding the use of coal versus natural gas?
    \end{enumerate}

\section*{Temperature}
  \textbf{Guiding questions to be addressed in your report:}
    \begin{enumerate}
    \item Describe the extent of surface warming so far, and its distribution as a function of latitude, for land versus ocean and for winter versus summer.
    \item Explain the latitudinal dependence; what are the implications for high-latitude countries? For Greenland melting and thus for other regions?
    \item Explain how future warming is evaluated from present-day data in addition to the use of climate models.
    \item If a doubling of CO$_2$ from 280 to 560 ppm leads to say, a 3\,K warming, what warming do you expect at an equivalent CO$_2$ mixing ratio of 450 ppm? Explain why the warming by 2020 was only about 1.1\,K.
    \item Should we be concerned about global warming given that it is said to have stopped recently for 15 yr? Will it stop again?
    \end{enumerate}

\section*{Sea-level}
  \textbf{Guiding questions to be addressed in your report:}
    \begin{enumerate}
    \item How and why is sea level changing along the Atlantic coast of the United States?
    \item How is past GMSL rise calculated, and what are the uncertainties (touch on spatial patterns and observational coverage of ocean warming now versus two hundred years ago, wind changes, isostatic response to glacial melting, coastal erosion, gravitational fingerprint effects, and so on).
    \item What is the projected GMSL rise by year 2100 based on RCP8.5? List the main processes responsible for the expected rise, quantify the contribution of each process, and describe the sources of uncertainty in each.
    \item Discuss the possible effect of anthropogenic global warming on storm surges and GMSL rise, and their interaction in specific regional sea level rise events.
    \item Discuss the expected effect of the projected GMSL rise by 2100 on New York City and Bangladesh. What are your recommendations in each case?
    \end{enumerate}

\section*{Acidification}
  \textbf{Guiding questions to be addressed in your report:}
    \begin{enumerate}
    \item What is ocean acidification, and why is it an important concern?
    \item What is the current state of ocean acidification relative to 1850, and what is expected by the year 2100 in a business-as-usual scenario?
    \item What is the ocean carbonate system?
    \item How and why does the concentration of the carbonate ion CO$_3^{2-}$ affect the dissolution of calcium carbonate?
    \item Based on the approximate solution to the carbonate system, how does the addition of CO$_2$ to the atmosphere change the dissolved CO$_2$, the bicarbonate ion, the carbonate ion, and thus calcium carbonate dissolution? Discuss the buffer effect.
    \item What are the economic consequences of ocean acidification? (Research this beyond what was presented in class.)
    \item What are the implications of ocean acidification for the utility of solar geoengineering?
    \end{enumerate}

\section*{Ocean-circulation}
  \textbf{Guiding questions to be addressed in your report:}
    \begin{enumerate}
    \item What is the AMOC, and why is it important?
    \item What controls its strength? Why might it change in a future climate? What are the implications of such a change?
    \item Discuss climate tipping points and explain why the AMOC may be an example.
    \item What are the consequences of such climate tipping points, in general, in terms of restoring climate to its preindustrial state once CO$_2$ is reduced back to preindustrial levels?
    \item What are the consequences of climate tipping points in terms of our ability to predict the occurrence of abrupt climate change and in terms of our ability to gradually adjust to changes?
    \end{enumerate}

\section*{Clouds}
  \textbf{Guiding questions to be addressed in your report:}
    \begin{enumerate}
    \item Why are clouds important in the current climate? Why are they important in the context of global warming?
    \item How do clouds form and dissipate, and what are the uncertainty sources in representing them in climate models?
    \item What are the mechanisms by which clouds can play a role in affecting the response to increased concentration of CO$_2$?
    \item What are the uncertainty sources in identifying and quantifying each of the mechanisms by which clouds affect the response to increased CO$_2$?
    \item Should we not worry about global warming because of the large uncertainty due to clouds?
    \end{enumerate}

\section*{Hurricanes}
  \textbf{Guiding questions to be addressed in your report:}
    \begin{enumerate}
    \item Have hurricanes (as quantified by the PDI, number of hurricanes per year, and the fraction of major hurricane winds) already been getting stronger?
    \item What are the factors affecting hurricane strength?
    \item How can the maximum possible (potential) hurricane strength be estimated from the SST\@? Why is it only a {\textit{potential}} maximum?
    \item Will hurricanes get stronger in an RCP8.5-like global warming scenario? If so, by how much? (Consider both hurricane wind velocity and destructiveness.)
    \item What are the sources of uncertainty?
    \end{enumerate}

\section*{Sea-ice}
  \textbf{Guiding questions to be addressed in your report:}
    \begin{enumerate}
    \item  Describe the changes in Arctic sea ice over the past decades.
    \item Why is the Arctic sea ice cover important?
    \item Why do these changes happen? Discuss all the positive and negative feedback mechanisms you can list and their possible roles in the observed changes.
    \item Can climate change be detected? Can the observed trend be attributed to anthropogenic global warming? Explain the difference between detection and attribution in this case.
    \item What are the sources of uncertainty in ACC detection? In the attribution to greenhouse forcing?
    \end{enumerate}

\section*{Ice-sheets}
  \textbf{Guiding questions to be addressed in your report:}
    \begin{enumerate}
    \item What are the climate and socioeconomic effects of a possible Greenland ice sheet melt in a warmer future climate?
    \item List, explain and compare different processes that may lead to future changes to the ice sheets in Greenland and Antarctica.
    \item List processes that may lead to a {\textit{rapid}} ice sheet collapse/retreat.
    \item Discuss factors that may make this retreat irreversible on a timescale of hundreds and possibly thousands of years. Which portions of which ice sheets may be affected?
    \item Suggest what elements of Greenland and Antarctica should be monitored to identify early warning signals.
    \end{enumerate}

\section*{Mountain-glaciers}
  \textbf{Guiding questions to be addressed in your report:}
    \begin{enumerate}
    \item Why are mountain glaciers important?
    \item How did they change over the past 200 yr?
    \item What are possible reasons for changes seen in mountain glaciers over different periods during the past 200 yr?
    \item What are the timescales of observed and expected changes?
    \item Investigate and discuss projections for the next 20, 50, and 100~yr for different glaciated mountain regions at both low and high latitudes, not including Greenland and Antarctica.
    \item Discuss the socioeconomic consequences of the projected further retreat of mountain glaciers, focusing on a couple of key regions as examples.
    \end{enumerate}

\section*{Droughts}
  \textbf{Guiding questions to be addressed in your report:}
    \begin{enumerate}
    \item Why are droughts important to understand and predict?
    \item Why and how can they change due to natural climate variability? Due to anthropogenic climate change?
    \item How do precipitation deficit and surplus events affect soil moisture? What are the roles of the duration and magnitude of such deficit and surplus?
    \item How and based on what data can we determine if a given drought may be attributable to anthropogenic climate change? What are the uncertainties in such a determination?
    \item Discuss projected drought conditions for the Southwest United States and for the Sahel and their expected socioeconomic consequences.
    \end{enumerate}

\section*{Floods}
\textbf{Guiding questions to be addressed in your report:}
    \begin{enumerate}
    \item What are the types of floods? Why are floods important to understand and predict?
    \item Why and how can each type change due to natural climate variability? Due to anthropogenic climate change? Due to other anthropogenic effects?
    \item Define extreme precipitation events. What are their socioeconomic consequences? How and why are they expected to change? 
    \item Why are changes to the probability of rare precipitation events difficult to estimate using observations? Why are rare flood events difficult to attribute to anthropogenic climate change?
    \item What can we say about the projections of changes to global-scale time-mean precipitation patterns in a warmer climate? What about regional changes?
    \end{enumerate}

\section*{Heat-waves}
  \textbf{Guiding questions to be addressed in your report:}
    \begin{enumerate}
    \item Why do heat waves occur in the present climate?
    \item What are the different factors playing a role in their mechanism?
    \item How is the character of these events expected to change in a warmer future climate? Why?
    \item What are the implications for human health, agriculture, and other activities? In your answer, consider separately regions that are currently classified as having temperate climates versus regions that are currently deserts or semi-arid climates, according to the {\textit{K\"oppen climate classification}}.
    \item Given that heat waves are defined as extreme and, therefore, rare events, what are the expected difficulties in attributing such individual events to anthropogenic climate change? Why is it still easier to attribute heat waves to ACC than drought events?
    \end{enumerate}

\section*{Forest-fires}
  \textbf{Guiding questions to be addressed in your report:}
    \begin{enumerate}
    \item Describe trends in burnt forest area in the western United States and Canada for the past decades.
    \item Discuss all {\textit{non-climate}}-related anthropogenic influences on forest fires and anthropogenic factors that affect fire damage.
    \item Explain the different ways warming and increases and decreases in precipitation rates may lead to changes in fire trends.
    \item Discuss the uncertainty sources involved in making future projections of fires on regional and global scales; discuss two specific regions as examples.
    \item Recommend action apart from an effort to reduce greenhouse emissions.
    \end{enumerate}

\end{document}
